
\documentclass{article}
\usepackage{colortbl}
\usepackage{makecell}
\usepackage{multirow}
\usepackage{supertabular}

\begin{document}

\newcounter{utterance}

\centering \large Interaction Transcript for game `wordle', experiment `medium\_frequency\_words\_no\_clue\_no\_critic', episode 11 with reflexion{-}Qwen3.5{-}4B.
\vspace{24pt}

{ \footnotesize  \setcounter{utterance}{1}
\setlength{\tabcolsep}{0pt}
\begin{supertabular}{c@{$\;$}|p{.15\linewidth}@{}p{.15\linewidth}p{.15\linewidth}p{.15\linewidth}p{.15\linewidth}p{.15\linewidth}}
    \# & \multicolumn{2}{c}{Player} && \multicolumn{2}{c}{Game Master} \\
    \hline

    \theutterance \stepcounter{utterance}  
    & & & \multicolumn{4}{p{0.6\linewidth}}{
        \cellcolor[rgb]{0.9,0.9,0.9}{
            \makecell[{{p{\linewidth}}}]{
                \texttt{\tiny{[P1$\langle$GM]}}
                \texttt{You have attempted this same task before and did not succeed. These are the plans you wrote after those attempts.} \\
\texttt{Trial \#1: New plan: The last attempt was aborted because the word "area" contains the letter 'a' twice, which violates the standard Wordle rule that each letter must be unique. Additionally, the explanation mentioned moving 'a' to a different position without explicitly stating the new word's structure. Next time, I will ensure every guess is a five{-}letter word with five distinct letters. I will also explicitly confirm in the explanation that the new word uses the known valid letter ('a') exactly once and replaces the invalid letters ('p', 'l', 'e') with other common letters, ensuring the response is both logically sound and strictly formatted.} \\
\\ 
\\ 
\texttt{You are a language wizard who likes to guess words by using the given rules.} \\
\\ 
\texttt{Welcome to Wordle! You have six attempts to guess the target word, a valid English word of five lowercase letters (a{-}z). Please use the tags "explanation:" and "guess:" to provide a concise explanation for each guess.} \\
\\ 
\texttt{For instance, if your guess is "apple", your response should be} \\
\texttt{explanation: this is a common five{-}letter English word, and I am starting my guess with this word.} \\
\texttt{guess: apple} \\
\\ 
\texttt{After each guess, your answer will be validated, and you will receive feedback indicating which letters are correct (green), which letters are correct but in the wrong position (yellow), and which letters are incorrect (red). This feedback can be useful in determining which letters to include or exclude in your next guess.} \\
\\ 
\texttt{For example, the feedback for "apple" might be:} \\
\texttt{guess\_feedback: a<yellow> p<yellow> p<green> l<yellow> e<red>} \\
\\ 
\texttt{The explanation should contain details about how the guess\_feedback is used to arrive at a new guess.} \\
\\ 
\texttt{Do not generate any other text.} \\
\\ 
\texttt{Let's begin with your first guess.} \\
            }
        }
    }
     \\ \\

    \theutterance \stepcounter{utterance}  
    & \multicolumn{4}{p{0.6\linewidth}}{
        \cellcolor[rgb]{0.561,0.737,0.561}{
            \makecell[{{p{\linewidth}}}]{
                \texttt{\tiny{[P1$\rangle$GM]}}
                \texttt{explanation: I will start with a common five{-}letter word with five distinct letters to gather information about the target word's letters and positions.} \\
\texttt{guess: crate} \\
            }
        }
    }
    & & \\ \\

    \theutterance \stepcounter{utterance}  
    & & & \multicolumn{2}{p{0.3\linewidth}}{
        \cellcolor[rgb]{0.9,0.9,0.9}{
            \makecell[{{p{\linewidth}}}]{
                \texttt{\tiny{[GM$|$GM]}}
                \texttt{attempts = 1} \\
\texttt{target\_word = faint} \\
\texttt{guess = crate} \\
\texttt{guess\_feedback = c<red> r<red> a<yellow> t<yellow> e<red>} \\
            }
        }
    }
    & & \\ \\

    \theutterance \stepcounter{utterance}  
    & & & \multicolumn{4}{p{0.6\linewidth}}{
        \cellcolor[rgb]{0.9,0.9,0.9}{
            \makecell[{{p{\linewidth}}}]{
                \texttt{\tiny{[P1$\langle$GM]}}
                \texttt{guess\_feedback: c<red> r<red> a<yellow> t<yellow> e<red>} \\
\\ 
\texttt{Provide your response only in this format.} \\
\texttt{explanation: details} \\
\texttt{guess: word} \\
            }
        }
    }
     \\ \\

    \theutterance \stepcounter{utterance}  
    & \multicolumn{4}{p{0.6\linewidth}}{
        \cellcolor[rgb]{0.561,0.737,0.561}{
            \makecell[{{p{\linewidth}}}]{
                \texttt{\tiny{[P1$\rangle$GM]}}
                \texttt{explanation: The feedback 'a<yellow> t<yellow>' tells me that 'a' and 't' are in the word but in the wrong positions. Since 'c', 'r', and 'e' are red, they are not in the word. I need a new word containing 'a' and 't' where 'a' is not in the second position and 't' is not in the fourth position.} \\
\texttt{guess: state} \\
            }
        }
    }
    & & \\ \\

    \theutterance \stepcounter{utterance}  
    & & & \multicolumn{2}{p{0.3\linewidth}}{
        \cellcolor[rgb]{0.9,0.9,0.9}{
            \makecell[{{p{\linewidth}}}]{
                \texttt{\tiny{[GM$|$GM]}}
                \texttt{attempts = 2} \\
\texttt{target\_word = faint} \\
\texttt{guess = state} \\
\texttt{guess\_feedback = s<red> t<yellow> a<yellow> t<red> e<red>} \\
            }
        }
    }
    & & \\ \\

    \theutterance \stepcounter{utterance}  
    & & & \multicolumn{4}{p{0.6\linewidth}}{
        \cellcolor[rgb]{0.9,0.9,0.9}{
            \makecell[{{p{\linewidth}}}]{
                \texttt{\tiny{[P1$\langle$GM]}}
                \texttt{guess\_feedback: s<red> t<yellow> a<yellow> t<red> e<red>} \\
\\ 
\texttt{Provide your response only in this format.} \\
\texttt{explanation: details} \\
\texttt{guess: word} \\
            }
        }
    }
     \\ \\

    \theutterance \stepcounter{utterance}  
    & \multicolumn{4}{p{0.6\linewidth}}{
        \cellcolor[rgb]{0.561,0.737,0.561}{
            \makecell[{{p{\linewidth}}}]{
                \texttt{\tiny{[P1$\rangle$GM]}}
                \texttt{explanation: The feedback 't<yellow> a<yellow>' confirms 't' and 'a' are in the word but in the wrong positions. 's' and 'e' are red, so they are excluded. I need a word with 'a' and 't' where 'a' is not in the third position and 't' is not in the second position.} \\
\texttt{guess: taste} \\
            }
        }
    }
    & & \\ \\

    \theutterance \stepcounter{utterance}  
    & & & \multicolumn{2}{p{0.3\linewidth}}{
        \cellcolor[rgb]{0.9,0.9,0.9}{
            \makecell[{{p{\linewidth}}}]{
                \texttt{\tiny{[GM$|$GM]}}
                \texttt{attempts = 3} \\
\texttt{target\_word = faint} \\
\texttt{guess = taste} \\
\texttt{guess\_feedback = t<yellow> a<green> s<red> t<red> e<red>} \\
            }
        }
    }
    & & \\ \\

    \theutterance \stepcounter{utterance}  
    & & & \multicolumn{4}{p{0.6\linewidth}}{
        \cellcolor[rgb]{0.9,0.9,0.9}{
            \makecell[{{p{\linewidth}}}]{
                \texttt{\tiny{[P1$\langle$GM]}}
                \texttt{guess\_feedback: t<yellow> a<green> s<red> t<red> e<red>} \\
\\ 
\texttt{Provide your response only in this format.} \\
\texttt{explanation: details} \\
\texttt{guess: word} \\
            }
        }
    }
     \\ \\

    \theutterance \stepcounter{utterance}  
    & \multicolumn{4}{p{0.6\linewidth}}{
        \cellcolor[rgb]{0.561,0.737,0.561}{
            \makecell[{{p{\linewidth}}}]{
                \texttt{\tiny{[P1$\rangle$GM]}}
                \texttt{explanation: The feedback 'a<green>' confirms 'a' is in the second position. 't' is yellow, meaning it is in the word but not in the first position (since it's green in the second spot of "taste", wait, the feedback says t<yellow> a<green>. In "taste", 't' is first, 'a' is second. The feedback says 'a' is green (correct position) and 't' is yellow (correct letter, wrong position). This is a contradiction because if 'a' is green in the second position, 't' cannot be yellow in the first position if the word is "state" or similar. Let's re{-}evaluate. The target word has 'a' in the second position. 't' is in the word but not in the first position. 's', 'e' are out. I need a word with 'a' in the second spot and 't' somewhere else.} \\
\texttt{guess: atone} \\
            }
        }
    }
    & & \\ \\

    \theutterance \stepcounter{utterance}  
    & & & \multicolumn{2}{p{0.3\linewidth}}{
        \cellcolor[rgb]{0.9,0.9,0.9}{
            \makecell[{{p{\linewidth}}}]{
                \texttt{\tiny{[GM$|$GM]}}
                \texttt{attempts = 4} \\
\texttt{target\_word = faint} \\
\texttt{guess = atone} \\
\texttt{guess\_feedback = a<yellow> t<yellow> o<red> n<green> e<red>} \\
            }
        }
    }
    & & \\ \\

    \theutterance \stepcounter{utterance}  
    & & & \multicolumn{4}{p{0.6\linewidth}}{
        \cellcolor[rgb]{0.9,0.9,0.9}{
            \makecell[{{p{\linewidth}}}]{
                \texttt{\tiny{[P1$\langle$GM]}}
                \texttt{guess\_feedback: a<yellow> t<yellow> o<red> n<green> e<red>} \\
\\ 
\texttt{Provide your response only in this format.} \\
\texttt{explanation: details} \\
\texttt{guess: word} \\
            }
        }
    }
     \\ \\

    \theutterance \stepcounter{utterance}  
    & \multicolumn{4}{p{0.6\linewidth}}{
        \cellcolor[rgb]{0.561,0.737,0.561}{
            \makecell[{{p{\linewidth}}}]{
                \texttt{\tiny{[P1$\rangle$GM]}}
                \texttt{explanation: The feedback 'n<green>' confirms 'n' is in the fourth position. 'a' and 't' are yellow, meaning they are in the word but not in the first and second positions respectively. 'o' and 'e' are red, so they are excluded. I need a word with 'n' in the fourth spot, 'a' and 't' elsewhere, and no 'o' or 'e'.} \\
\texttt{guess: tanka} \\
            }
        }
    }
    & & \\ \\

    \theutterance \stepcounter{utterance}  
    & & & \multicolumn{2}{p{0.3\linewidth}}{
        \cellcolor[rgb]{0.9,0.9,0.9}{
            \makecell[{{p{\linewidth}}}]{
                \texttt{\tiny{[GM$|$GM]}}
                \texttt{attempts = 5} \\
\texttt{target\_word = faint} \\
\texttt{guess = tanka} \\
\texttt{guess\_feedback = t<yellow> a<green> n<yellow> k<red> a<red>} \\
            }
        }
    }
    & & \\ \\

    \theutterance \stepcounter{utterance}  
    & & & \multicolumn{4}{p{0.6\linewidth}}{
        \cellcolor[rgb]{0.9,0.9,0.9}{
            \makecell[{{p{\linewidth}}}]{
                \texttt{\tiny{[P1$\langle$GM]}}
                \texttt{guess\_feedback: t<yellow> a<green> n<yellow> k<red> a<red>} \\
\\ 
\texttt{Provide your response only in this format.} \\
\texttt{explanation: details} \\
\texttt{guess: word} \\
            }
        }
    }
     \\ \\

    \theutterance \stepcounter{utterance}  
    & \multicolumn{4}{p{0.6\linewidth}}{
        \cellcolor[rgb]{0.561,0.737,0.561}{
            \makecell[{{p{\linewidth}}}]{
                \texttt{\tiny{[P1$\rangle$GM]}}
                \texttt{explanation: The feedback 'a<green>' confirms 'a' is in the third position. 'n' is yellow, so it is in the word but not in the fourth position. 't' is yellow, so it is in the word but not in the first position. 'k' is red, so it is excluded. I need a word with 'a' in the third spot, 'n' and 't' elsewhere, and no 'k'.} \\
\texttt{guess: tango} \\
            }
        }
    }
    & & \\ \\

    \theutterance \stepcounter{utterance}  
    & & & \multicolumn{2}{p{0.3\linewidth}}{
        \cellcolor[rgb]{0.9,0.9,0.9}{
            \makecell[{{p{\linewidth}}}]{
                \texttt{\tiny{[GM$|$GM]}}
                \texttt{attempts = 6} \\
\texttt{target\_word = faint} \\
\texttt{guess = tango} \\
\texttt{guess\_feedback = t<yellow> a<green> n<yellow> g<red> o<red>} \\
            }
        }
    }
    & & \\ \\

    \theutterance \stepcounter{utterance}  
    & & & \multicolumn{2}{p{0.3\linewidth}}{
        \cellcolor[rgb]{0.9,0.9,0.9}{
            \makecell[{{p{\linewidth}}}]{
                \texttt{\tiny{[GM$|$GM]}}
                \texttt{game\_result = LOSS} \\
            }
        }
    }
    & & \\ \\

    \theutterance \stepcounter{utterance}  
    & & \multicolumn{4}{p{0.6\linewidth}}{
        \cellcolor[rgb]{0.95,0.95,0.95}{
            \makecell[{{p{\linewidth}}}]{
                \texttt{\tiny{[P1$\leftrightarrow$P]}}
                \texttt{Post{-}episode Reflexion (not sent to GM):} \\
\texttt{New plan: The previous attempt failed because I misinterpreted the feedback logic, specifically confusing the meaning of "green" and "yellow" in relation to the target word's positions. I assumed 'a' was green in the second position based on the word "tango" (where 'a' is actually in the third position), leading to a contradiction. I need to strictly adhere to the feedback: if a letter is green, it is in the correct position; if yellow, it is in the word but in the wrong position. My new strategy will be to carefully track the exact position of every green letter and the exclusion of all red letters. For the next guess, I will construct a word that places 'a' in the third position (since 'tango' feedback showed 'a' as red, wait, let me re{-}read the last feedback: "tango" gave "t<yellow> a<green> n<yellow> k<red> a<red>". Wait, "tango" has 'a' in the 2nd and 4th positions. The feedback says 'a' is green. This implies the target word has 'a' in the 2nd position. But my previous analysis said 'a' was green in the 3rd spot? No, the feedback for "tango" was t<yellow> a<green> n<yellow> k<red> a<red>. The second letter of "tango" is 'a'. So 'a' is green, meaning it is in the 2nd position. The 4th letter is 'o', which is red. The 3rd letter is 'n', which is yellow. The 1st is 't' (yellow), 5th is 'o' (red). So the target has 'a' in pos 2, 'n' in pos 3 (since 'n' is yellow in "tango" at pos 3? No, "tango" is t{-}a{-}n{-}g{-}o. Pos 3 is 'n'. Feedback says n<yellow>. So 'n' is in the word but not pos 3. Pos 1 't' is yellow, so 't' is not pos 1. Pos 2 'a' is green, so 'a' IS pos 2. Pos 4 'g' is red. Pos 5 'o' is red. So I have: \_ a \_ \_ \_. 'n' is in the word but not pos 3. 't' is in the word but not pos 1. I need a word with 'a' in pos 2, 'n' and 't' elsewhere, and no 'g' or 'o'. My previous guess "tango" failed because I thought 'a' was green in the 3rd spot, but it was the 2nd spot. I need} \\
            }
        }
    }
    & \\ \\

\end{supertabular}
}

\end{document}
